% Style commun des lettres de motivation — Ismaila AMADOU DEH
\documentclass[11pt,a4paper]{article}
\usepackage[utf8]{inputenc}
\usepackage[T1]{fontenc}
\usepackage[english,french]{babel}
\usepackage{lmodern}
\usepackage[margin=2.2cm,top=1.8cm]{geometry}
\usepackage{xcolor}
\usepackage{parskip}
\usepackage[hidelinks]{hyperref}
\definecolor{accent}{HTML}{1F4E79}
\definecolor{gris}{HTML}{6B675F}
\pagestyle{empty}
% Coordonnées privées (e-mail, téléphone) : fichier local non versionné
\newcommand{\contact}{}
\IfFileExists{coordonnees.tex}{\input{coordonnees.tex}}{}
\newcommand{\entete}{%
  {\LARGE\bfseries\color{accent} Ismaila AMADOU DEH}\\[2pt]
  {\small\color{gris} Étudiant M2 Physique Fondamentale -- parcours Théories Quantiques, Université de Tours\\
  \contact\href{https://www.linkedin.com/in/ismaila-amadou-deh}{linkedin.com/in/ismaila-amadou-deh} \;·\;
  \href{https://github.com/IsmailaDEH/Ismaila}{github.com/IsmailaDEH/Ismaila}}\\[-4pt]
  {\color{accent}\rule{\linewidth}{0.6pt}}}
\newcommand{\destinataire}[2]{\null\hfill\begin{minipage}[t]{0.45\linewidth}#1\\[6pt]Tours, le \today\end{minipage}\par\medskip
  \textbf{Objet :} #2\par}
\newcommand{\signature}{\par\bigskip\hfill Ismaila AMADOU DEH}

\begin{document}
\entete
\destinataire{Thales Research \& Technology\\Service Recrutement\\Palaiseau}{Candidature au stage « Ingénieur en caractérisation de circuits photoniques intégrés »}

Madame, Monsieur,

Étudiant en Master~2 Physique Fondamentale, parcours Théories Quantiques, à l'Université de Tours, je vous adresse ma candidature pour ce stage. La photonique intégrée est au carrefour de ce qui m'attire : l'optique expérimentale, la mesure de précision et les technologies quantiques.

Durant mon Master~1, j'ai réalisé de nombreux travaux pratiques d'optique et de mesure : \textbf{spectroscopie laser He-Ne et sodium}, couplage spin-orbite, résonance paramagnétique électronique et magnétométrie optique par centres NV. J'y ai appris à aligner un montage, à mener une mesure soignée et à interpréter un spectre en tenant compte des incertitudes.

Je suis également à l'aise avec le traitement de données et l'automatisation. Je développe en Python des chaînes d'analyse complètes : ajustement de modèles physiques, extraction de paramètres avec barres d'erreur, génération automatique des figures. Par exemple, dans mon projet de spectroscopie d'un ion piégé, je relie l'asymétrie des bandes latérales rouge et bleue à la température de l'ion. Ce savoir-faire se transpose directement à la caractérisation de circuits photoniques (pertes, résonances, couplage).

Rigoureux et curieux, je souhaite apprendre les techniques de caractérisation sur banc que vous utilisez et contribuer à l'automatisation des mesures.

Je serais heureux de vous rencontrer lors d'un entretien. Veuillez agréer, Madame, Monsieur, l'expression de mes salutations distinguées.
\signature
\end{document}
